\begin{examplebox}
\par
\begin{minipage}{\linewidth}
\textbf{\textcolor{CExample}{例 1｜标准椭圆：求极线并核对两切点}}\par\smallskip
\textbf{题目：}椭圆 $x^2/9+y^2/4=1$，$P(3,1)$，求极线并说明是否为切点弦。

\begin{center}
\begin{tikzpicture}[x=1.55cm,y=1.55cm,>=Stealth,font=\small]
 \path[use as bounding box] (-3.5,-2.4) rectangle (4,2.7);
 \draw[->,gray] (-3.3,0)--(3.7,0) node[right] {$x$};
 \draw[->,gray] (0,-2.2)--(0,2.35) node[above] {$y$};
 \draw[CStatement,very thick] (0,0) ellipse (3 and 2);
 \draw[orange!80!black,thick] (3,-.6)--(3,2.1);
 \draw[orange!80!black,thick] (.8,2.1)--(3.8,.6);
 \draw[teal!80!black,very thick] (1.25,2.33333)--(3.4,-.53333);
 \fill (3,1) circle(1.8pt) node[above right] {$P(3,1)$};
 \fill (3,0) circle(1.8pt) node[below left] {$T_1(3,0)$};
 \fill (1.8,1.6) circle(1.8pt);
 \draw[gray] (1.8,1.6)--(.9,2.15);
 \node at (.7,2.5) {$T_2(\frac95,\frac85)$};
 \node[teal!80!black] at (-1.45,-2.25) {切点弦 $4x+3y=12$};
\end{tikzpicture}

\end{center}
图中一条切线竖直，故直接写极线比先设两条斜率更稳妥。

\end{minipage}
\par
\textbf{解：}代入椭圆公式：
\[
 \frac{3x}{9}+\frac{y}{4}=1
 \quad\Longleftrightarrow\quad\boxed{4x+3y=12}.
\]
\textbf{不要在这里停下：}要称为切点弦，还要确认这条线与椭圆有两个不同实交点。
\par\medskip
由直线得 $x=(12-3y)/4$，代入椭圆：
\begin{align*}
 &\frac{(12-3y)^2}{16\cdot9}+\frac{y^2}{4}=1,\\
 &144-72y+9y^2+36y^2=144,\\
 &9y(5y-8)=0,
\end{align*}
所以 $y=0$ 或 $y=8/5$，相应 $x=3$ 或 $x=9/5$。
两个切点分别是
\[
 T_1(3,0),\qquad T_2\left(\frac95,\frac85\right).
\]
验算：在 $T_1$ 处的切线为 $x=3$，含 $P(3,1)$；
在 $T_2$ 处的切线为 $x/5+2y/5=1$，即 $x+2y=5$，也含 $P(3,1)$。
因此 $4x+3y=12$ 确为切点弦所在直线。

\medskip
\textbf{\textcolor{CExample}{例 2｜抛物线：曲线上一点的切线}}\par\smallskip
\textbf{题目：}抛物线 $y^2=8x$ 上一点 $P(2,4)$，求切线。

\textbf{解：}$2p=8$，故 $p=4$。先验算 $4^2=8\cdot2$，点确实在曲线上。
代入 $y_0y=p(x+x_0)$：
\[
 4y=4(x+2)\quad\Longrightarrow\quad\boxed{y=x+2}.
\]
检查相切：代入曲线得到
\[
 (x+2)^2=8x\quad\Longleftrightarrow\quad(x-2)^2=0.
\]
$x=2$ 是重根，切点正是 $(2,4)$。

\medskip
\textbf{\textcolor{CExample}{例 3｜含 $xy$ 的一般二次方程：代换与几何核对}}\par\smallskip
\textbf{题目：}$\Gamma:x^2+xy+y^2-2x-2y+1=0$，$P(2,0)$，求极线和切点。

\textbf{解：}对照一般式，
$A=B=C=1,\ D=E=-2,\ F=1$。逐项代入：
\begin{align*}
 0={}&2x+\frac{2y}{2}+0-(x+2)-y+1\\
  ={}&x-1.
\end{align*}
故 $\boxed{x=1}$。
将 $x=1$ 代入 $\Gamma$，得到
\[
 1+y+y^2-2-2y+1=0
 \quad\Longleftrightarrow\quad y(y-1)=0.
\]
切点为 $T_1(1,0)$、$T_2(1,1)$。
分别求切线：在 $T_1$ 处是 $y=0$；在 $T_2$ 处是 $x+y=2$。
两条切线都经过 $P(2,0)$，再次确认 $x=1$ 是切点弦。

\medskip
\par
\begin{minipage}{\linewidth}
\textbf{\textcolor{CExample}{例 4｜双曲线：有极线但只有一个实切点}}\par\smallskip
\textbf{题目：}双曲线 $x^2-y^2=1$，$P(1,1)$。能否把极线称为切点弦？

\begin{center}
\begin{tikzpicture}[x=1.65cm,y=1.65cm,>=Stealth,font=\small]
 \path[use as bounding box] (-2.6,-2) rectangle (4.2,2.2);
 \draw[->,gray] (-2.3,0)--(3.3,0) node[right] {$x$};
 \draw[->,gray] (0,-1.8)--(0,1.9) node[above] {$y$};
 \draw[CStatement,very thick,domain=-1.5:1.5,samples=70] plot ({sqrt(1+\x*\x)},\x);
 \draw[CStatement,very thick,domain=-1.5:1.5,samples=70] plot ({-sqrt(1+\x*\x)},\x);
 \draw[gray,dashed] (-1.7,-1.7)--(1.8,1.8);
 \node[gray,rotate=45,fill=white,inner sep=1pt] at (-1,-1) {$y=x$};
 \draw[teal!80!black,very thick] (-.65,-1.65)--(2.65,1.65);
 \draw[orange!80!black,thick] (1,-1.65)--(1,1.8);
 \fill (1,1) circle(1.8pt) node[above left] {$P(1,1)$};
 \fill (1,0) circle(1.8pt) node[below right] {$T(1,0)$};
 \node[teal!80!black,anchor=west] at (2.1,.6) {极线 $x-y=1$};
 \node[orange!80!black,anchor=west] at (1.3,2) {切线 $x=1$};
\end{tikzpicture}

\end{center}
极线平行于渐近线 $y=x$，只交曲线于 $T$；橙色竖线才是从 $P$ 引出的切线。

\end{minipage}
\par
\textbf{解：}极线方程为
\[
 \boxed{x-y=1}.
\]
因为曲线上 $x^2-y^2=(x-y)(x+y)=1$，同时又有 $x-y=1$，故必有 $x+y=1$。
解得唯一实交点 $T(1,0)$。
双曲线在 $T$ 处的切线是 $x=1$，的确经过 $P(1,1)$。
\textbf{结论：}这条极线存在，但只有一个不同的实切点，\textbf{不能称为两切点的连线}。

\medskip
\textbf{\textcolor{CExample}{例 5｜最小反例：椭圆内点与中心}}\par\smallskip
对椭圆 $x^2/9+y^2/4=1$：
\begin{itemize}[leftmargin=*,itemsep=0.35em]
 \item $P(1,0)$ 的极线是 $x=9$。椭圆上 $|x|\le3$，所以它没有实交点，也没有实切点弦。
 \item $P(0,0)$ 的代换结果是 $0=1$，不能写作一条直线，更不能叫切点弦。
\end{itemize}
\end{examplebox}
